\chapter{Entorno y preparación de datos}

\noindent El Capítulo~\ref{chap:propuesta} define el estado del modelo MARL como $s_t=(m_t,c_t,h_t)$: información de mercado, estado de cartera y señal supervisada. Este capítulo explica el origen de los dos bloques que dependen de información externa. Los registros actuales de los mercados se normalizan para construir $m_t$, mientras que sus series históricas permiten obtener las variables y la señal supervisada $h_t$. El bloque de cartera $c_t$ no procede de una fuente externa: el simulador lo inicializa y actualiza a partir de las decisiones ejecutadas.

El capítulo identifica las fuentes de información disponibles, explica cómo se normalizan y validan los registros para conservar su procedencia, su fecha de captura y los atributos que identifican cada artículo, y presenta las variables y el procedimiento temporal con el que se forman los conjuntos de entrenamiento, validación y prueba. Esta preparación permite que los modelos utilicen únicamente la información disponible en cada fecha y que sus resultados se evalúen sin incorporar información futura.

\section{Fuentes de datos y relación con el estado}

El Capítulo~\ref{chap:caso-estudio} describe los mercados secundarios que delimitan el caso de estudio. Este apartado concreta el uso de cada fuente dentro del flujo de datos. El flujo principal comienza con SteamDT y continúa con la consulta de Steam Community Market y \acroref{BUFF}{BUFF} para los artículos seleccionados:

\begin{itemize}
  \item \textbf{SteamDT} genera una lista inicial de artículos candidatos \cite{steamDT}. Esta lista reduce el número de artículos que deben consultarse en las fuentes de precio.
  \item \textbf{Steam Community Market} proporciona datos actuales e históricos de precio, órdenes de compra y volumen de intercambios \cite{steamMarketFAQ}. Su histórico se utiliza para construir el dataset temporal.
  \item \textbf{\acroref{BUFF}{BUFF}} proporciona el precio actual de compra y las órdenes de compra con los que se contrasta una operación de compraventa entre plataformas \cite{buffMarketplace}. No se utiliza como fuente histórica principal del entrenamiento.
  \item \textbf{Reconocimiento óptico de caracteres (\acroref{OCR}{OCR}) e importación manual:} son dos vías complementarias al flujo principal. El OCR extrae texto y cifras de una captura de pantalla. La importación manual permite cargar una tabla o un fichero preparado previamente. Ambas vías se emplean cuando los datos no están disponibles mediante una consulta automatizada.
\end{itemize}

Cada registro almacenado conserva la fuente, la fecha de captura, la moneda original y los atributos que identifican el artículo. La siguiente sección describe cómo se normalizan y validan estos campos. Los procedimientos técnicos de captura se presentan en el Capítulo 6.

La Tabla~\ref{tab:fuentes-estado-marl} resume la trazabilidad desde cada fuente hasta el modelo de decisión. Una misma fuente puede contribuir al bloque de mercado actual y, si proporciona una serie temporal suficiente, a la construcción de la señal supervisada.

\begin{table}[H]
\centering
\caption{Relación entre las fuentes de datos y el estado del modelo MARL.}
\label{tab:fuentes-estado-marl}
\begin{tabularx}{\textwidth}{@{}p{0.22\textwidth}p{0.33\textwidth}X@{}}
\toprule
\textbf{Origen} & \textbf{Datos conservados} & \textbf{Destino en $s_t$} \\
\midrule
SteamDT & Lista inicial de artículos candidatos. & Delimita los candidatos que pueden originar un bloque $m_t$; no produce por sí solo una señal supervisada. \\
Steam Community Market & Precios actuales e históricos, órdenes de compra y volumen de intercambios. & Aporta campos de $m_t$ y, mediante su histórico, variables para calcular $h_t$. \\
BUFF & Precio actual de compra y órdenes de compra. & Completa los precios y la disponibilidad de $m_t$ para la ruta de compraventa entre mercados. \\
OCR e importación manual & Precios, moneda y atributos extraídos de capturas o ficheros. & Sustituyen o complementan campos de $m_t$ y solo contribuyen al histórico si superan la validación descrita en la Sección~5.2. \\
Simulador de cartera & Saldo inicial y consecuencias de compras, ventas y bloqueos. & Construye y actualiza $c_t$; no es una fuente de datos de mercado. \\
\bottomrule
\end{tabularx}
\end{table}

\section{Normalización y validación de los datos}

Una fuente puede expresar los mismos datos con nombres, monedas, fechas o formatos distintos. Para evitar que esas diferencias se interpreten como diferencias de mercado, cada captura se transforma en un \textbf{registro de mercado}: una fila de dato que junta un artículo, la fuente de la que procede, la fecha de captura de datos y los valores publicados en ese momento. Este registro es la unidad que se utiliza después para construir variables temporales.

La normalización conserva el nombre y los atributos que identifican el artículo, como su acabado, nivel de desgaste, \textit{float} o StatTrak, y los asocia a una misma identificación interna. También permite comparar precios publicados en monedas distintas. Por ejemplo, si Steam Community Market muestra el precio de un artículo en euros y \acroref{BUFF}{BUFF} lo muestra en yuanes, el sistema convierte ambos precios a una moneda de comparación. Al mismo tiempo, conserva el precio y la moneda publicados por cada fuente para poder comprobar posteriormente de dónde procede cada valor. Esta conversión solo permite comparar precios. Las comisiones, el beneficio neto y la rentabilidad se calculan después mediante las reglas económicas definidas en el Capítulo~\ref{chap:propuesta}.

La fecha es igualmente necesaria. Los registros se ordenan por fecha de captura para que una variable calculada en un día solo emplee datos disponibles hasta ese día. Si una fuente proporciona históricos, cada punto histórico mantiene su propia fecha en lugar de recibir la fecha de descarga. De este modo, el conjunto temporal distingue cuándo se publicó un dato de cuándo el sistema lo incorporó.

Antes de almacenar un registro, el sistema comprueba que el artículo pueda identificarse y que los campos necesarios para su tipo de dato estén presentes. Por ejemplo, un precio histórico sin fecha o sin valor numérico no forma parte de la serie temporal. Asimismo, un mismo artículo, fuente, fecha y tipo de dato, como precio, orden de compra o volumen de intercambios, no se incorpora dos veces. Los datos obtenidos mediante \acroref{OCR}{OCR} deben superar además sus controles específicos de lectura y confianza, explicados en el Capítulo 6. Los detalles de captura, persistencia y deduplicación se describen también en ese capítulo.

\section{Variables y conjunto temporal}

Los registros normalizados no se emplean directamente para entrenar los modelos. A partir de ellos se construyen \textbf{variables de entrada}, es decir, valores numéricos que resumen la situación de un artículo en una fecha determinada. Todas las variables de una fila se calculan únicamente con los datos disponibles hasta esa fecha. Por tanto, el resultado posterior de una posible operación no forma parte de la información utilizada para predecirla.

Las variables de mercado se agrupan en los siguientes bloques:

\begin{itemize}
  \item \textbf{Precios y comparación entre mercados:} precio de venta publicado en Steam Community Market y en \acroref{BUFF}{BUFF}, precio de las órdenes de compra de \acroref{BUFF}{BUFF} y diferencias entre esos precios antes y después de aplicar las comisiones correspondientes.
  \item \textbf{Actividad del mercado:} volumen de intercambios registrado en Steam Community Market y número de ofertas publicadas en \acroref{BUFF}{BUFF}. Estos valores indican cuánta actividad se ha producido alrededor del artículo.
  \item \textbf{Evolución temporal:} cambios y retornos de precio a 1, 3 y 7 días, así como la media y la desviación de los últimos 7 días. También se incorporan el día de la semana, el mes y un indicador de fin de semana.
\end{itemize}

Los valores derivados de la comparación de precios en la fecha inicial, como el beneficio neto y el ROI calculados con los precios disponibles en esa fecha, pueden utilizarse como variables de entrada. El modelo supervisado transforma esas variables en la probabilidad que forma $h_t$. No se incorporan al conjunto de datos las variables de cartera, porque el simulador las crea y actualiza durante cada episodio, tal como se describe en el Capítulo~\ref{chap:propuesta}.

Para formar el conjunto temporal, se toma un artículo $i$ y una fecha $t$. El sistema calcula sus variables con los registros disponibles hasta $t$ y busca después el primer registro futuro de precio dentro del intervalo comprendido entre $t+8$ y $t+15$ días. El primer límite representa el periodo de bloqueo de intercambio utilizado en la simulación. Los siete días adicionales permiten conservar el ejemplo cuando la fuente no publica un registro exactamente en el octavo día.

La etiqueta $y_{i,t}$ vale 1 cuando la operación simulada con ese registro futuro produce un beneficio neto no negativo y un ROI de al menos el 10\%. En cualquier otro caso, vale 0. Si $d$ representa la fecha del primer registro futuro encontrado, $B_{i,d}$ el beneficio neto calculado y $R_{i,d}$ el ROI, la etiqueta se define como:

\[
y_{i,t} = 1
\Longleftrightarrow
B_{i,d} \geq 0
\quad \land \quad
R_{i,d} \geq 0{,}10.
\]

Esta etiqueta no predice una simple subida de precio. Indica si, bajo las reglas económicas del Capítulo~\ref{chap:propuesta}, una compra en la fecha $t$ habría alcanzado el criterio económico definido al llegar la fecha $d$.

El conjunto actual separa los datos por fecha para evitar que los modelos conozcan información futura:

\begin{itemize}
  \item \textbf{Entrenamiento:} desde el 2 de julio de 2025 hasta el 16 de diciembre de 2025.
  \item \textbf{Validación:} desde el 1 de enero de 2026 hasta el 13 de febrero de 2026.
  \item \textbf{Prueba:} desde el 1 de marzo de 2026 hasta el 30 de junio de 2026.
\end{itemize}

Los quince días anteriores a cada frontera se excluyen del conjunto. Esta franja de purga impide que un registro de entrenamiento o validación utilice como resultado un precio perteneciente al periodo siguiente. Con este procedimiento, cada conjunto conserva únicamente la información disponible en su periodo y permite evaluar los modelos sin incorporar datos futuros. Cuando el entorno MARL recorre un episodio, consulta la probabilidad producida con estas variables temporales para completar $h_t$; los precios y atributos de la fecha del paso completan $m_t$ y el simulador actualiza $c_t$.
